% TOP_PROBLEM: {"id": 2303034, "title": "Research Problems in Function Theory — Problem 3.34", "category_id": 8, "category": {"id": 8, "name": "analysis", "display_name": "Analysis", "description": "Limits, continuity, calculus, and function theory.", "slug": "analysis", "order_index": 8, "created_at": "2026-07-31T15:26:25.671Z"}, "status": "open"}
% FIRST_POSTED: null
% ENTRY_KIND: statement_only
% Imported from: batch13.tex; UnsolvedMath ID 2303034
\documentclass[11pt]{book}
\usepackage{fontspec}
\setmainfont{Latin Modern Roman}
\setsansfont{Latin Modern Sans}
\usepackage{amsmath,amssymb,mathtools}
\usepackage{unicode-math}
\setmathfont{Latin Modern Math}
\usepackage{xurl}
\usepackage[hidelinks]{hyperref}
\newcommand{\sref}[2]{\hyperlink{source:#1:#2}{[S#2]}}
\newcommand{\cataloguescope}[1]{\paragraph{Mathematical question.}#1}
\begin{document}
% UnsolvedMath ID 2303034; source catalogue code AMR-022-3034
\setcounter{section}{2303033}
\section{Research Problems in Function Theory — Problem 3.34}\label{2303034}
{\small\sffamily UnsolvedMath ID 2303034 · Analysis}
\subsection{Definitions and mathematical statement}

\paragraph{Shared notation.}$\mathbb N=\{1,2,\ldots\}$, and $\mathbb Z,\mathbb Q,\mathbb R,\mathbb C$ denote the integers, rationals, reals and complex numbers. Any other conventions are specified locally.


Fix $n\ge2$ and $0<\alpha\le\pi$, interpreting $\alpha$ as the full opening angle of a circular cone. Let $\Omega\subset\mathbb R^n$ be a bounded connected open set such that for every $y\in\partial\Omega$ there are a unit vector $e_y$ and height $h_y>0$ with
\[\{y+te:0<t<h_y,\ |e|=1,\ e\cdot e_y>\cos(\alpha/2)\}\subset\Omega.
\]
The height is allowed to depend on the boundary point; no Lipschitz boundary or uniform positive lower height is imposed. A positive superharmonic function $V$ is a lower semicontinuous function $V:\Omega\to(0,\infty]$, not identically infinite, satisfying the super-mean inequality on every ball compactly contained in $\Omega$. For $p>0$, the statement $V\in L^p(\Omega)$ means $\int_\Omega V(x)^p\,dx<\infty$.
\paragraph{Statement recorded in the supplied source.}
Does there exist an exponent $p=p_n(\alpha)>0$, independent of the particular domain and function, such that every positive superharmonic $V$ on every such $\Omega$ belongs to $L^p(\Omega)$? If so, characterize the possible exponents in terms of the cone angle (with dimension fixed). \sref{2303034}{2}
\subsection{Short English statement}
Does having an interior truncated cone of a fixed angle at each boundary point force a common positive integrability exponent for all positive superharmonic functions, and how does that exponent depend on the angle?
\subsection{Sources}
\begin{description}
\item[\hypertarget{source:2303034:1}{S1}] ULAM.AI and the UnsolvedMath Contributors, \emph{UnsolvedMath: A Curated Collection of Open Mathematics Problems}, record 2303034 (AMR-022-3034). CC BY 4.0. \url{https://huggingface.co/datasets/ulamai/UnsolvedMath}.
\item[\hypertarget{source:2303034:2}{S2}] Walter K. Hayman and Edward F. Lingham, \emph{Research Problems in Function Theory}, arXiv:1809.07200, 2018. Problem 3.34, its complete Update 3.34, and the relevant chapter notation. \url{https://arxiv.org/pdf/1809.07200}.
\end{description}
\label{end:2303034}
\end{document}
